\scanpage{28}
\mspage{26}
% archive sheet number: 27
% hand: Grothendieck
% sheet_kind: text
% element: text | upper, lines 1–4 | 5.7 Comme groupe algébrique
5.7. Comme groupe algébrique abstrait \unsure{cependant}, $G$ est déterminé par la connaissance de $\Gamma\simeq\pi/U$, et de sa représentation par autom. de \deleted{\unsure{$U/N\otimes\mathbb Z$}} du vectoriel \deleted{\unsure{$F$}$E$}
% element: display | upper, line 5 | U/N tensor Q_l
\[
U/N\otimes_{\mathbb Z_\ell}\mathbb Q_\ell,
\]
% element: text | upper, line 6 | qui est un quotient
qui est un quotient du vectoriel
% element: display | upper-middle, lines 7–8 | U/[U,U]^(ell) tensor Q_l
\[
\underbrace{U/[U,U]^{(\ell)}}_{\substack{\text{\unsure{composante|complété}}\\\ell\text{-primaire}}}\otimes_{\mathbb Z_\ell}\mathbb Q_\ell.
\]
% element: text | middle, lines 9–11 | Lorsque pi=I
Lorsque $\pi=I\simeq\inserted{\operatorname{Gal}(\overline K/K)\simeq}\mathbb Z_\ell(1)$ comme au n° 2, \inserted{(cas str.\textsuperscript{t} local)} on sait, \deleted{que} puisque $U$ est d'indice fini, qu'on a \inserted{et} un isomorphisme canonique
% element: display | middle, line 12 | U/[U,U]^(ell) isom Z_l(1)
\[
U/[U,U]^{(\ell)}\simeq\mathbb Z_\ell(1)
\tag*{(5.7.1)}
\]
% element: text | middle, lines 13–14 | isom canonique
(isom canonique), \deleted{\unsure{d'où on a une}} une représentation naturelle de $\pi/U=\Gamma$ ``dans $\mathbb Z_\ell(1)$'', i.e.
% element: display | middle-lower, line 15 | pi/U=Gamma to Z_l^*
\[
\pi/U=\Gamma\longrightarrow\mathbb Z_\ell^*
\]
% element: text | lower, lines 16–17 | qui est connue
qui est connue \uline{pour être triviale} \inserted{(3.15)}. Donc en fait on aura un isom. \inserted{canonique} déduit de (5.5.3)
% element: display | lower, line 18 | G isom G^0 times Gamma
\[
G\simeq G^\circ\times\Gamma\quad\deleted{\unsure{\simeq E}}
\tag*{(5.7.2)}
\]
% element: text | lower, lines 19–22 | où G^0 est soit nul
\deleted{et $G^\circ$ est le groupe} \uline{où $G^\circ$ est soit nul, soit} \deleted{\unsure{égal}} \uline{canoniquement isomorphe au} \deleted{\unsure{$\mathbb Q_\ell(1)$}} \uline{groupe algébrique défini par le vectoriel} $\mathbb Q_\ell(1)$. On a en particulier
% element: display | bottom, lines 23–24 | cases for G(Q_l)
\[
\left\{\begin{aligned}
G(\mathbb Q_\ell)&\simeq\Gamma=\pi/U)&&\text{si }G\text{ fini, i.e. }U\text{ opère trivialement}\\
G(\mathbb Q_\ell)&\simeq\mathbb Q_\ell(1)\times\Gamma&&\text{sinon.}
\end{aligned}\right.
\tag*{(5.7.3)}
\]
